%%=============================================================================
%% Conclusie
%%=============================================================================

\chapter{Conclusie}%
\label{ch:conclusie}

Deze bachelorproef onderzocht hoe computervisie kan bijdragen aan tijdige indicaties van lokale drukte tijdens openluchtevenementen zoals de Gentse Feesten. De literatuurstudie toont dat lokale drukte niet uitsluitend door het aantal aanwezigen wordt bepaald. Ruimte, doorstroming, zichtbaarheid, de duur van een verhoogde bezetting en de samenwerking tussen diensten zijn eveneens bepalend. Een automatische meting kan menselijke observatie daarom ondersteunen, maar niet vervangen.

Voor dichte menigtes is density estimation een geschikte aanvulling op persoonsdetectie. Waar objectdetectie voor iedere persoon een afzonderlijk begrenzingskader moet vinden, voorspelt een density-estimationmodel een ruimtelijke verdeling en een totale telling. Daardoor is de methode beter bestand tegen overlap en schaalverschillen, al blijft zij afhankelijk van representatieve trainingsdata. CSRNet werd gekozen als eerste baseline omdat de architectuur voor dicht bezette scènes is ontworpen en op ShanghaiTech goed is gedocumenteerd.

De ontwikkelde PoC realiseert een reproduceerbare keten van puntannotaties naar density maps, training, full-image-validatie en visualisatie. Op ShanghaiTech Part B werd de officiële trainingsset met een vaste seed opgesplitst in 360 trainingsbeelden en 40 validatiebeelden. De officiële testset werd niet gebruikt voor modelselectie. In de finale GPU-run van 100 epochs ligt het checkpoint met de laagste validatie-MAE op epoch 39. Het behaalt een validatie-MAE van \(6{,}556\) personen, een RMSE van \(9{,}218\) personen en een bijna neutrale bias van \(-0{,}304\) personen. Tegenover een gecontroleerde run van één epoch met dezelfde configuratie dalen MAE en RMSE met respectievelijk \(76{,}5\%\) en \(77{,}4\%\). De volledige training duurt \(58{,}19\) minuten. Voor een beeld van \(1024 \times 768\) pixels bedraagt de gemiddelde tijd voor uitsluitend de voorwaartse modelberekening op de RTX 3050 Ti ongeveer \(107\) ms, of 9,3 FPS.

De centrale onderzoeksvraag kan bijgevolg als volgt worden beantwoord: computervisie kan met density estimation een bruikbare aanvullende indicatie geven van het geschatte aantal personen en van relatieve druktezones in een camerabeeld. Voor de Gentse Feesten is dat potentieel waardevol als één invoer voor een operator. De huidige PoC bewijst echter nog niet dat het systeem voldoende nauwkeurig, robuust of snel is voor operationeel gebruik. Met ongeveer 9,3 model-doorlopen per seconde voldoet de gemeten GPU-inferentie niet aan een eventuele doelwaarde van 15 beelden per seconde. Op de Raspberry Pi 5-CPU duurt één model-doorloop gemiddeld \(10{,}95\) s en wordt slechts \(0{,}091\) FPS bereikt. Daarbij trad thermische throttling op. Voor een beoordeling van real-timewerking moeten ook beeldopname, voorbewerking, visualisatie en communicatie worden gemeten.

Een cruciale beperking is het onderscheid tussen een density map en fysieke personendichtheid. De PoC werkt in beeldcoördinaten en kan zonder camerakalibratie, grondvlaktransformatie en afgebakende observatiezone geen waarde in personen per vierkante meter berekenen. Zij mag daarom geen zelfstandige alarmdrempel hanteren. Ook gebruikt ShanghaiTech Part B andere camerastandpunten en omgevingsomstandigheden dan de Gentse Feesten. Nachtelijke belichting, weersinvloeden, lokale architectuur en veranderende bezoekersstromen kunnen tot domeinverschuiving leiden. Representatieve, rechtmatig verkregen beelden uit de doelcontext zijn nodig om die generaliseerbaarheid te beoordelen.

Verder onderzoek moet de finale configuratie met meerdere seeds herhalen en CSRNet vergelijken met minstens één alternatief model, zoals MCNN. Ook de cropgrootte vraagt een gecontroleerde vergelijking. Uitsneden van bijvoorbeeld \(256\), \(384\), \(512\) en \(768\) pixels kunnen met dezelfde datasplit, hetzelfde aantal epochs en dezelfde overige hyperparameters worden getraind en op dezelfde volledige validatiebeelden worden beoordeeld. Zo kan worden nagegaan of extra beeldcontext de MAE en RMSE verbetert en hoeveel extra trainingstijd en GPU-geheugen dit kost. Vervolgens moeten de latentie van de volledige verwerkingsketen, het geheugenverbruik en de nauwkeurigheid per camera worden gemeten. Voor fysieke dichtheidsdrempels zijn per camera een kalibratie, een grondvlakprojectie en een validatieprocedure nodig. Ten slotte moet een reële uitrol worden aangevuld met gegevensbeschermingsmaatregelen, waaronder een duidelijke rechtsgrond, dataminimalisatie, beperkte bewaartermijnen en een beoordeling van de gegevensbescherming. Onder die voorwaarden kan de huidige PoC uitgroeien tot een geloofwaardige bouwsteen voor beslissingsondersteuning bij crowd management.
